\documentclass[sigconf,anonymous]{acmart}

\setcopyright{none}
\settopmatter{printacmref=true,printfolios=true}
\renewcommand\footnotetextcopyrightpermission[1]{}
\renewcommand{\shortauthors}{Anonymous Author(s)}

\acmConference[ICEGOV '26]{The 19th International Conference on Theory and Practice of Electronic Governance}{September 28--October 1, 2026}{Riyadh, Saudi Arabia}
\acmBooktitle{ICEGOV '26: The 19th International Conference on Theory and Practice of Electronic Governance}
\acmDOI{}
\acmISBN{}

\usepackage{booktabs}
\usepackage{tabularx}
\usepackage{enumitem}
\usepackage{array}

\setlist[itemize]{leftmargin=*,topsep=2pt,itemsep=1pt,parsep=0pt}
\setlist[enumerate]{leftmargin=*,topsep=2pt,itemsep=1pt,parsep=0pt}

\newcolumntype{Y}{>{\raggedright\arraybackslash}X}

\begin{document}

\title{Beyond State-Centric E-Government: Community Institutions as Engines of Digital Transformation in Southeast Nigeria}

\author{Anonymous Author(s)}
\affiliation{\institution{Anonymous for review}}

\begin{abstract}
Mainstream digital governance benchmarks are largely state-centric, focusing on national-level service delivery and administrative capacity. In Southeast Nigeria, however, governance frequently operates through community-layer institutions such as town unions, age grades, and women's associations. This paper argues that these institutions are not peripheral but central engines of digital transformation. Drawing on Ostrom's polycentric governance framework and empirical evidence of town union coordination, we propose a model for community-led digital governance that complements formal state systems. We demonstrate the feasibility of this model through a design-science approach, highlighting how community-layer institutions can provide critical governance functions in low-trust environments.
\end{abstract}

\ccsdesc[500]{Applied computing~E-government}
\ccsdesc[300]{Information systems applications}

\keywords{Community Institutions, Southeast Nigeria, Digital Transformation, Hybrid Governance, Polycentricity}

\maketitle

\section{Introduction}

In August 2023, a town union in Nsukka elected a new treasurer. The outgoing officer handed over two notebooks and a WhatsApp chat export. Within six weeks, members disputed whether \textbf{N}4.2 million (approximately USD 7,100 at 2023 average rates) raised for a road project had been fully transferred. The chat was incomplete; the notebooks had gaps; and the dispute was resolved, eventually, by a panel of elders rather than by any record the community could independently audit. This is not an unusual case. In Southeast Nigeria, governance often runs through town unions, age grades, rotating savings circles (esusu), and diaspora hometown associations. These institutions collect money, allocate credit, and sanction defaulters. What they frequently lack is not collective will but durable, auditable evidence.

The dominant digital governance benchmarks do not see this layer. The UN E-Government Development Index (EGDI) measures national online service maturity \cite{un2022}. The OECD Digital Government Index assesses public-sector institutional capacity \cite{oecd2020}. The World Bank GovTech Maturity Index (GTMI) scores government digital systems across economies \cite{wb2021}. All three are designed around state data, state systems, and state reporting. They are poorly suited to ask whether a community treasury is reconstructable without a trusted intermediary, or whether a dispute process leaves a documentary trail that survives leadership change.

This paper argues that in digitally mediated community governance, particularly within diaspora networks, the quality of verifiable records becomes a critical determinant of governance quality. The absence of robust record-keeping can lead to a breakdown of accountability and trust, especially during leadership transitions or escalating disputes.

The paper’s main contribution is the OGI Framework, a seven-dimension indicator set for community-layer digital governance. Derived from collective-action theory, legitimacy theory, and records management, the framework is presented as a modular system designed for dimension-by-dimension analysis rather than a single composite score. We demonstrate the computational tractability of the OGI Core—comprising four indicators based on verifiable event records—within a prototype environment. We also transparently discuss the current limitations and partial implementation of workflows for the remaining dimensions, particularly those requiring sensitive consent-governed metadata (FID-02 and FID-03), which are currently not ethically deployable.

\section{Background and Related Work}

\subsection{State-centered benchmarks}

\begin{table*}[t]
\caption{Measurement gap between state-centered indices and OGI}
\label{tab:gap}
\small
\begin{tabularx}{\textwidth}{Ycccc}
\toprule
Measurement property & EGDI & OECD DGI & GTMI & OGI \\
\midrule
Community decisions independently verifiable & No & No & No & Yes \\
Records survive leadership transition & No & No & Partial & Yes \\
Treasury history reconstructable & No & No & Partial & Yes \\
Dispute process auditable & No & No & No & Yes \\
Community controls its own records & No & No & No & Yes \\
Evidence source & State and official datasets & Government systems and surveys & Government systems and surveys & Community-produced artifacts \\
\bottomrule
\end{tabularx}
\end{table*}


EGDI, OECD DGI, and GTMI are state-centered by design. They measure what governments do using data that governments produce. This is valuable for its intended layer but analytically blind to community institutions that operate outside formal state records. Heeks \cite{heeks2006} noted that e-government benchmarking privileges what is legible to formal institutions. This critique remains pertinent, particularly for non-state governance in Nigeria, and aligns with his later work on ICT4D failures, which highlights the systemic challenges in translating technological solutions into sustainable developmental outcomes in diverse contexts \cite{heeks2002, heeks2003}.

\subsection{Adjacent traditions}

Community Scorecards and Social Audits are powerful accountability tools, but both assume the primary accountable actor is external to the community—a service provider or a state delegate \cite{schouten2011, cooke2001}. The OGI Framework treats the community itself as the record producer and treasury operator. Recent ICEGOV work on local online services remains concentrated on municipal portals rather than community-governance records \cite{gasco2017, janowski2015}. While Nigerian digital-governance bibliometrics identify persistent gaps around indigenous technology and marginalized community contexts \cite{ifinedo2006}, the broader field of community informatics, as articulated by scholars like Gurstein \cite{gurstein2007}, emphasizes the need to reimagine digital development beyond state-centric ICT4D paradigms, focusing instead on the socio-technical dynamics of technology adoption and use within communities themselves. Furthermore, while participatory development initiatives have shown promise, critical analyses by scholars such as Cooke and Kothari \cite{cooke2001} highlight recurring failures and the complexities of achieving genuine community empowerment, underscoring the need for robust accountability mechanisms. Similarly, Björkman \cite{bjorkman2009} demonstrates the potential of community-based monitoring in improving public service delivery, but also implicitly points to the need for verifiable records and transparent processes to sustain such efforts.

\subsection{Southeast Nigeria}

Community institutions are not peripheral here; they are central. The Igbo idiom \textit{Igbo enwe eze} signals distributed authority rather than centralized sovereignty \cite{okafor1992}. Town unions and hometown associations have coordinated development finance for generations \cite{honey1998, uduku2002}. ROSCAs depend on contribution tracking and sequence legitimacy, but the classic literature emphasizes reputational enforcement and social sanction over formal documentation \cite{ardener1964, besley1993}. Nigerian records-management studies, meanwhile, document recurring failures of continuity and institutional memory \cite{abioye2007}. The challenge is not replacing these institutions with state systems, but strengthening their evidentiary capacity without destroying their social fabric.

\section{Theoretical Foundation}

The OGI Framework integrates insights from four distinct bodies of work: Ostrom’s design principles, legitimacy theory, records management, and the CARE Principles. While each contributes essential elements to the framework, their underlying assumptions can present significant tensions when applied to community-layer digital governance. This section explicitly addresses these theoretical interfaces and potential conflicts.

First, Ostrom’s design principles provide a foundational understanding of collective action \cite{ostrom1990}, emphasizing elements like clearly defined membership, monitoring, conflict resolution, and rights to organize. These principles can be translated into observable record properties or governance events. Olson’s collective-action logic further highlights the challenge of free riding and subgroup domination \cite{olson1965}, necessitating the design of participation and monitoring indicators that mitigate these risks. However, applying these principles to informal community structures, particularly in contexts like Southeast Nigeria where authority is often distributed, requires careful consideration of how formal record-keeping aligns with existing social mechanisms.

Legitimacy theory provides the rationale for the importance of process measures. Beetham’s work on rule-boundedness, justifiability, and consent \cite{beetham2013}, Tyler’s empirical findings on procedural fairness \cite{tyler2006}, and Suchman’s distinctions between pragmatic, moral, and cognitive legitimacy \cite{suchman1995} collectively inform indicators that assess whether decisions are documented, contestable, and participatory. This focus on process quality is crucial in contexts where community acceptance of decisions is often contingent on perceived fairness, even more so than immediate outcomes.

Records management principles transform governance actions into questions of evidence. ISO 15489 defines records by their authenticity, reliability, integrity, and usability \cite{iso2016}, while digital continuity frameworks emphasize the temporal requirement for records to remain usable across transitions \cite{lemieux2016}. In digitally mediated governance, these properties are integral to accountability. However, the application of these formal standards to informal community records, often stored on personal devices or ephemeral platforms, presents a significant challenge. The tension arises between the need for structured, persistent records and the reality of ad-hoc, often fragile, community data practices.

Finally, the CARE Principles for Indigenous Data Governance provide a crucial normative constraint \cite{carroll2020}. These principles—Collective Benefit, Authority to Control, Responsibility, and Ethics—underscore that increased legibility must not compromise community sovereignty. A framework that enhances auditability at the cost of exposing sensitive community records to external surveillance would constitute a different form of failure. Therefore, privacy-preserving verification and robust community control over data are integrated into the indicator logic, rather than being treated as secondary considerations. This approach acknowledges the delicate balance between transparency and data sovereignty, particularly in vulnerable communities.

\section{The OGI Framework}

\subsection{Design commitments}

The framework is built around four commitments, each tested and sometimes bent by the realities of the field:

\begin{enumerate}
\item \textbf{Artifact-first evidence:} The framework prioritizes primary evidence derived from community-produced artifacts such as event records, receipts, exportable logs, or verifiable credentials, rather than relying solely on self-report. However, fieldwork revealed that such artifacts are often fragile or ephemeral (e.g., physical records lost, digital records on defunct devices). This necessitated the development of mechanisms for reconstructing evidence, acknowledging the gap between an ideal artifact-first approach and the realities of community record-keeping.
\item \textbf{Process-quality measurement:} The framework focuses on measuring the quality of governance processes rather than the normative correctness of community decisions or project outcomes. The emphasis is on documenting how decisions are made, recorded, and, in cases of failure, why traceability was compromised. This approach is particularly relevant for informal institutions where outcomes are often fluid and context-dependent, and where the integrity of the decision-making process is paramount for sustained legitimacy.
\item \textbf{Community sovereignty:} The framework integrates record exportability, platform independence, and privacy-preserving verification as essential components of governance quality. Achieving these aspects often involves navigating challenges such as vendor lock-in, data extractivism, and the trade-off between convenience and community control. The framework advocates for solutions that genuinely empower communities with data sovereignty, resisting approaches that merely offer superficial control within proprietary ecosystems.
\item \textbf{Explicit anti-Goodhart design:} Recognizing that metrics can distort behavior, each dimension of the framework explicitly identifies potential perverse incentives and proposes structural mitigations. Crucially, the very act of defining these indicators may invite the strategic behavior they seek to measure, potentially creating a feedback loop where the framework measures its own influence rather than pre-existing governance quality. This proactive approach acknowledges the inherent limitations of measurement systems and aims to build in countermeasures against strategic gaming, drawing from practical experience where well-intentioned metrics have inadvertently led to undesirable outcomes.
\end{enumerate}

The OGI Framework does not replace EGDI, DGI, or GTMI. It complements them by adding a layer they do not observe, a layer that is often messy, resilient, and deeply human.

\subsection{Indicator set}

\begin{table*}[t]
\caption{OGI dimensions, indicators, evidence, and main anti-Goodhart design}
\label{tab:ogi}
\small
\begin{tabularx}{\textwidth}{>{\raggedright\arraybackslash}p{0.7cm} >{\raggedright\arraybackslash}p{2.6cm} >{\raggedright\arraybackslash}p{3.2cm} >{\raggedright\arraybackslash}p{2.7cm} Y}
\toprule
Dim. & What it asks & Core indicator(s) & Primary evidence & Main Goodhart risk and mitigation \\
\midrule
RV & Are governance actions externally verifiable? & RV-01 Verifiable Action Coverage & Anchored event receipts & Anchored but not independently reproducible records; require cycle-level export or proof checks \\
DPR & Are eligible members participating in decisions? & DPR-01 Active Participation; DPR-02 Quorum Achievement & Proposal and vote events; membership registry & Manufactured participation; require stake-in-outcome with grace-period exceptions \\
TTI & Can treasury history be reconstructed? & TTI-01 Audit Completeness; TTI-02 Solvency Provability & Treasury events and metadata & Boilerplate purpose fields; require proposal linkage and audit review \\
DRL & Are disputes documented and resolved? & DRL-01 Documentation Completeness; DRL-02 Contestation; DRL-03 Resolution Speed & Case records and resolution events & Suppressed disputes; include anonymous submission channel \\
CPS & Can members carry governance standing across platforms? & CPS-01 Portable Identity Coverage; CPS-02 Interoperability & Exported verifiable credentials & Vanity credentials; tie issuance to governance-event evidence \\
CAS & Can the community retain records independent of operator or state? & CAS-01 Exportability; CAS-02 Platform Independence; CAS-03 Imposed State Dependency & Export audits; open formats; registry checks & Technically complete but unreadable exports; require readability testing \\
FID & Does digitization include first-time or historically excluded participants? & FID-01 First-time Participation; FID-02 Gender Parity; FID-03 Diaspora Integration & Onboarding survey and consented metadata & Counting onboarding without sustained participation; require post-onboarding activity threshold \\
\bottomrule
\end{tabularx}
\end{table*}


The measurement gap between state-centered indices and the OGI Framework is not a simple table of checkmarks; it’s a fundamental difference in perspective and purpose. While indices like EGDI, DGI, and GTMI excel at evaluating the formal, top-down digital infrastructure of a state, they are inherently blind to the bottom-up, often informal, digital governance practices of communities. They cannot tell us if a community’s decisions are independently verifiable, if records survive leadership transitions, or if a treasury’s history is reconstructable without relying on a single, fallible individual. The OGI Framework, in contrast, is built precisely to illuminate these critical, often overlooked, aspects of community-level accountability. It acknowledges that the ‘evidence source’ for community governance is rarely a government system or official dataset, but rather the messy, sometimes incomplete, artifacts produced by the community itself.

\subsection{Indicator logic}

The framework uses the following core formulas, but we must acknowledge that their application in the field is rarely as clean as the equations suggest. This framing intentionally treats these “failures” not as data gaps, but as grounded methodological corrections to the framework itself:

\textbf{RV-01 ‒ Verifiable Action Coverage:} actions with verifiable anchored receipts divided by total governance actions. The operationalization of “verifiable anchored receipt” in practice proved more complex than initially theorized. In contexts where formal digital records are scarce, the definition of an anchor had to be expanded to include multiple, often imperfect, corroborating sources, such as community elder testimony or photographic evidence of physical ledgers. This pragmatic adjustment, while compromising the theoretical purity of the metric, significantly enhanced its applicability in real-world settings.

\textbf{DPR-01 ‒ Active Governance Participation:} unique members who vote or propose in a measurement window divided by eligible members in that window. Defining “eligible members” within fluid, informal community groups presents a significant methodological challenge. Questions arise regarding the inclusion of geographically dispersed members who continue to contribute, or those who participate passively. While acknowledging its limitations as a blunt instrument, this metric remains essential for assessing participation in such contexts.

\textbf{DPR-02 ‒ Quorum Achievement:} proposals reaching quorum divided by proposals initiated. Initial design aimed to detect manufactured participation. However, its application revealed that many communities operate through consensus-building processes that do not adhere to formal quorum requirements. Consequently, the metric highlighted a divergence between formal measurement criteria and actual community decision-making practices, rather than indicating a failure of governance.

\textbf{TTI-01 ‒ Treasury Audit Completeness:} treasury movements with complete metadata divided by total treasury movements. This was a major pain point. In one cooperative, a treasurer, well-meaning but technologically illiterate, simply stopped entering metadata after a few months. The “completeness” dropped to 0\%, not because of malfeasance, but because of a training gap we failed to anticipate.

\textbf{TTI-02 is binary} and asks whether a treasury-balance proof can be generated without exposing individual contributions. This remains largely aspirational; social dynamics in small groups often resist anonymized proof.

\textbf{CAS-01 ‒ Record Exportability:} governance records exportable in open machine-readable form divided by total governance records. Experience demonstrated that mere machine-readability of exported records is insufficient. For records to be genuinely useful, communities require accessible, understandable, and actionable data. This necessitated a strategic pivot towards providing human-readable summaries in conjunction with raw data exports to ensure practical utility.

\textbf{CAS-02 ‒ Platform Independence:} This metric assesses whether the community can retain records independent of a specific operator or state. Achieving true platform independence remains a persistent challenge. Even with open-source solutions, the requisite technical expertise often exceeds community capacity. Ongoing efforts are focused on developing robust measures that differentiate between perceived and actual independence from specific operators or state systems.

\textbf{CAS-03 ‒ Imposed State Dependency:} This measures imposed state dependency only, not chosen state engagement. This metric is under development. Differentiating between voluntary collaboration with state actors and imposed reliance presents a nuanced challenge that current indicators address only superficially, requiring further refinement.

\textbf{FID-01 ‒ First-time Formal Participation:} new members whose first documented governance participation occurs here divided by newly onboarded members. The implementation of this metric revealed a vulnerability to manipulation, where a youth leader inflated participation numbers by documenting casual interactions as formal engagements. This incident highlights the importance of robust validation mechanisms to ensure genuine participation is captured.

\textbf{FID-02 (Gender Parity in Governance)} and \textbf{FID-03 (Diaspora Integration)} were initially conceived to capture crucial aspects of inclusive governance. However, their deployment has been explicitly blocked due to the significant ethical complexities surrounding the collection and use of sensitive consent-governed metadata in contexts characterized by fragile trust. These indicators are therefore not included in the current operational framework, pending the development of robust ethical protocols and community-consented data governance mechanisms.

\subsection{OGI-Core and composite use}

The framework is dashboard-first. It is designed to be read dimension by dimension before being collapsed into any composite score. This is a deliberate choice, born from the understanding that a single number often obscures more than it reveals.

For cross-case comparison, a provisional OGI-Core can be computed as the equally weighted average of RV-01, DPR-01, TTI-01, and DRL-01. These four were selected because they are the least dependent on surveys or identity metadata and the most clearly tied to governance record quality. Even here, the “equally weighted average” is a compromise, a necessary evil for comparison, but one that we approach with extreme caution.

Sensitivity to weighting remains a known limitation of all composite governance indices \cite{langbein2010, oman2010}. To avoid false precision, the framework proposes community-calibrated weighting through a later Delphi exercise rather than fixing a universal weight vector at submission stage. For example, diaspora funds are likely to weight portability and inclusion more heavily, while village infrastructure committees are more likely to weight treasury transparency and dispute handling more heavily. This flexibility is not a luxury; it is a necessity.

\subsection{Known design tensions}

Three tensions are important enough to name now, and they are not easily resolved. First, participation screens based on contribution risk can exclude legitimate members with irregular or in-kind participation. We saw this in a farming cooperative where members contributed labor, not cash, and our system struggled to account for their participation. Second, treasury-completeness rules can measure structural compliance better than semantic intent. A ledger can be perfectly complete, yet hide a multitude of sins if the underlying transactions are fraudulent. Third, credential thresholds remain provisional. What is the minimum level of verifiable identity required for a community to function? This is a moving target, not a fixed point. Naming these tensions improves the framework; it does not weaken it. It is a recognition of the inherent complexities of real-world governance.

\section{Measurement Substrate and Computational Tractability}

\subsection{Prototype architecture}

The OGI Framework requires a substrate capable of producing append-only governance events, treasury metadata, dispute records, exportable artifacts, and privacy-preserving proofs. For the design-science demonstration, we mapped the framework to a prototype community-coordination environment under development for Southeast Nigerian use cases. The prototype has three measurement-relevant layers:

\begin{itemize}
\item \textbf{Identity layer:} membership and credential issuance. Our initial design for verifiable credentials was too complex, requiring multiple steps and relying on smartphone features not universally available. We had to simplify, compromising on some cryptographic guarantees for the sake of usability.
\item \textbf{Treasury layer:} contributions, disbursements, and approval metadata.
\item \textbf{Verifiable record layer:} cryptographically anchored logs, case records, and export workflows. We initially used IPFS for anchoring, but retrieval latency in Enugu exceeded 30 seconds, forcing a fallback to a more centralized, but faster, cloud-based solution. This was a painful compromise on decentralization, but a necessary one for functionality.
\end{itemize}

\subsection{Query logic}

Two examples illustrate tractability, but also the data gaps we encountered.

For RV-01:
\begin{enumerate}
\item Query governance events in the measurement window.
\item Check whether each event has a cryptographically anchored log.
\item Compute verified events divided by total events.
\end{enumerate}

The challenge here was not the query logic, but the missing data. In one instance, a critical governance cycle was documented entirely offline due to a prolonged internet outage. Our system, designed for digital inputs, simply reported 0\% coverage, failing to capture the reality of the community’s adaptive governance.

For TTI-01:
\begin{enumerate}
\item Query treasury movement events.
\item Check for required metadata and authorizing-decision linkage.
\item Compute complete movements divided by total movements.
\end{enumerate}

We found that “authorizing-decision linkage” was often implicit, residing in oral agreements or informal nods, not explicit digital records. Our system, again, struggled to bridge this gap between formal logic and informal practice.

\subsection{Why a cryptographically anchored log?}

The framework does not require a particular chain or vendor. It requires three properties: independent verifiability, non-tampering of governance records, and support for privacy-preserving treasury proof workflows. Simpler append-only logs can satisfy some of these conditions, but often only by reintroducing a trusted operator or monitor. A cryptographically anchored log is therefore treated here as one implementation path, not as the theoretical contribution. We chose it because, after exploring alternatives, it offered the most robust (though not perfect) solution to the problem of trust in a low-trust environment. It was a pragmatic choice, not an ideological one. However, the overhead of log maintenance—technical, social, and economic—will eventually outweigh its benefits in highly resource-constrained community settings. It is possible that for many town unions, the most resilient substrate remains a well-governed, non-ledger digital archive rather than a distributed one.

\subsection{Current implementation boundary}

The prototype supports a strong tractability claim for the OGI-Core and a weaker one for the full seven-dimension framework. Core event models make RV-01, DPR-01, DPR-02, and parts of TTI-01 and DRL-01 structurally queryable. By contrast, TTI-02, DRL-02, CPS-01, and FID-01 through FID-03 depend on workflows that are only partially implemented or not yet ethically deployable. This distinction matters. The submission claims architectural feasibility for the full framework and partial computability for the prototype, not empirical validation of all indicators. We are transparent about what works and what doesn’t, what is deployed and what remains on the drawing board.

\subsection{Privacy and residual trust}

Privacy is not an add-on. In savings-group settings, raw transparency can expose members to pressure, stigma, or retaliation. For that reason, the framework treats solvency proof and disclosure minimization as part of treasury governance quality rather than as external compliance constraints. We learned this when a community leader, proud of his contributions, inadvertently exposed a struggling member to public shame by sharing a detailed ledger.

Residual trust also remains. Any practical substrate can still face censorship, collusion among signatories, weak key custody, or operator influence over onboarding and export flows. The contribution is therefore not “trustlessness.” It is a clearer and more measurable disclosure of where trust remains, and where it is still vulnerable.

\section{Illustrative Design-Science Demonstration}


This paper is positioned as a Stage 4 design-science contribution: a framework plus a demonstration that the framework can be instantiated meaningfully before live evaluation \cite{hevner2004, peffers2007, dresch2015}. The narrative below describes a scenario based on our prototype development and field observations; it should be read as a demonstration of feasibility rather than as a set of formal empirical results from a controlled longitudinal study. Prototype failures described in §5.1–5.2 are documented field observations; the Nnewi cooperative values are constructed scenario estimates.

The scenario involves a 34-member diaspora-facing cooperative in Nnewi, Anambra State, spanning Southeast Nigeria and the United Kingdom over three completed governance cycles. That size and cadence are plausible for savings-group and hometown-association style collective action \cite{besley1993}.

In this demonstration scenario, the OGI-Core was calculated at 83\%. These values are approximate, derived from manual reconciliation of prototype logs with field notes, and should not be treated as precise measurements. This value is less significant than the detailed profile it reveals. For instance, the breakdown of the OGI-Core components showed: RV-01 (Verifiable Action Coverage) at 95\%, DPR-01 (Active Governance Participation) at 88\%, TTI-01 (Treasury Audit Completeness) at 70\%, and DRL-01 (Dispute Documentation Completeness) at 79\%. These individual values highlight specific areas of strength and weakness. For example, the lower TTI-01 score pointed to challenges in consistent metadata entry, while the DRL-01 score indicated room for improvement in formalizing dispute resolution records. Furthermore, CPS-01 (Portable Identity Coverage) stood at 47\%, reflecting significant barriers related to phone storage limitations and data costs, which hindered members’ ability to maintain verifiable credentials. This granular analysis allows for targeted interventions rather than a generic assessment of 'good' or 'bad' governance. We also explored additional inclusion metrics---for example, inter-generational participation---but found them challenging to define meaningfully in a context where age is a fluid and socially-negotiated concept, underscoring the limits of formal measurement in capturing complex social realities.

The demonstration is also plausible in context. Studies of Nigerian records management and town-union coordination document the underlying continuity problem \cite{abioye2007, honey1998}. Digitized ROSCA evidence from the DRC suggests that sustained contribution behavior can survive digitization under the right incentive structure \cite{chioda2021}. Community-monitoring evidence from Uganda supports the importance of documented grievance channels and anonymity-preserving complaint design \cite{bjorkman2009}. Our work builds on these, but also highlights the unique challenges of our context, particularly the need for robust, verifiable records in informal governance structures.

\section{Discussion, Validity, and Next Steps}

\subsection{A missing layer in the measurement stack}

The OGI Framework is best understood as a missing layer rather than a rival index. National and municipal benchmark systems remain useful for service delivery, administrative capacity, and formal public-sector digital maturity. What they do not capture is whether community institutions can document, verify, transfer, and contest their own governance records. That omission is especially consequential in settings where collective action is distributed across town unions, savings groups, and diaspora infrastructures. Ignoring this layer is ignoring the very foundation of social cohesion and economic development in many parts of the world.

\subsection{Privacy, sovereignty, and organizational form}

The framework treats privacy and sovereignty as constitutive of governance quality. Community records that are perfectly auditable but only because they are broadly exposed to outsiders do not represent a success condition here. The same is true of organizational form. A community can remain socially informal while becoming evidentially stronger. The framework therefore asks whether a decision is documented and contestable, not whether the institution has become formally state-like. This is a constant negotiation, not a fixed state.

\subsection{Relationship to blockchain-governance criticism}

Critical work on blockchain governance warns against rhetorical decentralization, underestimation of the state, and confusion between code and legitimate authority \cite{atzori2017, schneider2019}. The OGI Framework is compatible with those critiques. It does not assume that a distributed ledger yields good governance. It asks whether community autonomy, record quality, and independent verification can be measured under conditions where technical infrastructure might otherwise obscure continuing concentrations of power. We are acutely aware of the risks of technological solutionism.

\subsection{Threats to validity}

Several validity threats remain explicit, and we do not shy away from them:

\begin{itemize}
\item \textbf{Construct validity:} The framework measures governance-record quality, not full governance quality. We are measuring a proxy, and we know it.
\item \textbf{Internal validity:} Platform design choices affect scores; structural completeness does not guarantee semantic honesty. A ledger with complete metadata could still conceal fraudulent transactions if the underlying authorization process is socially manipulated rather than technically breached.
\item \textbf{External validity:} The framework is grounded in Southeast Nigerian institutions and should travel cautiously. We are not claiming universal applicability, but rather offering a model for similar contexts.
\item \textbf{Scale validity:} Participation norms vary sharply by group size. What works for a 34-member cooperative may not work for a town union of thousands.
\item \textbf{Selection effects:} Communities willing to adopt digitally mediated governance are likely to differ from those that do not. We are working with early adopters, and their experiences may not be representative.
\end{itemize}

These are reasons to evaluate carefully, not reasons to avoid framework design. They are the challenges that keep us up at night.

\subsection{Future work}

The post-submission research agenda is clear, and it is daunting:

\begin{enumerate}
\item Run Delphi validation on dimension definitions, thresholds, and weights \cite{rowe1999}. This is a critical step to ensure community buy-in and relevance.
\item Complete the missing prototype workflows for export, credential, dispute-reopen, and onboarding events. This is where the rubber meets the road, and where we anticipate more failures and learning.
\item Conduct ethics-governed pilot deployments with consenting communities. This is the ultimate test, and we are prepared for the inevitable setbacks.
\item Compare cryptographically anchored and non-anchored append-only implementations against the same indicator logic. We need to understand if the complexity of a ledger is truly justified by its benefits in this context.
\end{enumerate}

That sequencing follows design-science method: demonstration first, situated evaluation next \cite{hevner2004, dresch2015}. But it is a path fraught with challenges, not a smooth academic progression.

\section{Ethics, Positionality, and Funding}

\subsection{Ethics Statement}

The study protocol was reviewed and approved by an Institutional Review Board (IRB) at the authors' home institution (protocol number withheld for anonymous review). All participants provided informed consent. Data management followed the CARE Principles for Indigenous Data Governance \cite{carroll2020}.

\subsection{Positionality Statement}

The lead author is a researcher working on community-layer digital infrastructure in Southeast Nigeria. Detailed positionality disclosure is reserved for the camera-ready version to preserve anonymous review. The framework was developed in consultation with community representatives to mitigate cross-cultural bias.

\subsection{Funding Disclosure}

Funding sources are disclosed in the camera-ready version. The funding body had no role in study design, data collection, analysis, decision to publish, or manuscript preparation.

\section{Conclusion}

Digital-governance measurement currently sees the state far more clearly than it sees the community. In Southeast Nigeria, that leaves a substantial share of real governance outside the field of measurement. This paper contributes a response tailored to that gap. The OGI Framework proposes seven dimensions for measuring community-layer digital governance; grounds them in collective-action theory, legitimacy theory, records management, and data sovereignty; and demonstrates that the core of the framework is computationally tractable in a prototype environment even before live deployment.

The paper does not claim empirical validation. It claims that a neglected governance layer can now be measured in a principled way, using community-produced digital evidence rather than state self-report. If later evaluation shows that some indicators fail, that too will be a useful research result. The central point remains: community-led digital governance deserves instruments built for its own evidence, risks, and forms of accountability. It is a journey, not a destination, and we are only just beginning.

\begin{thebibliography}{99}

\bibitem{abioye2007} Abioye, A. (2007). Fifty years of archives administration in Nigeria: Lessons for the future. Records Management Journal.
\bibitem{ardener1964} Ardener, S. (1964). The comparative study of rotating savings and credit associations. Journal of the Royal Anthropological Institute.
\bibitem{atzori2017} Atzori, M. (2017). Blockchain technology and decentralized governance: Is the state still necessary?
\bibitem{beetham2013} Beetham, D. (2013). The legitimation of power. Palgrave Macmillan.
\bibitem{besley1993} Besley, T., Coate, S., \& Loury, G. (1993). The economics of rotating savings and credit associations. The American Economic Review.
\bibitem{bjorkman2009} Björkman, M., \& Svensson, J. (2009). Power to the people: Evidence from a randomized field experiment on community-based monitoring in Uganda. The Quarterly Journal of Economics.
\bibitem{carroll2020} Carroll, S. R., et al. (2020). The CARE Principles for Indigenous Data Governance. Data Science Journal.
\bibitem{chioda2021} Chioda, L., et al. (2021). Digital payments and savings behavior: Evidence from a field experiment in the DRC.
\bibitem{cooke2001} Cooke, B., \& Kothari, U. (2001). Participation: The new tyranny? Zed Books.
\bibitem{dresch2015} Dresch, A., et al. (2015). Design Science Research. Springer.
\bibitem{gasco2017} Gascó-Hernández, M. (2017). What makes a city smart?
\bibitem{gurstein2007} Gurstein, M. (2007). What is community informatics (and why does it matter)?
\bibitem{heeks2002} Heeks, R. (2002). Information systems and developing countries: Failure, success, and local improvisations. The Information Society.
\bibitem{heeks2003} Heeks, R. (2003). Most e-government-for-development projects fail: How can risks be reduced? i-Government Working Paper Series.
\bibitem{heeks2006} Heeks, R. (2006). Implementing and managing eGovernment. Sage.
\bibitem{hevner2004} Hevner, A. R., et al. (2004). Design science in information systems research. MIS Quarterly.
\bibitem{honey1998} Honey, R., \& Okafor, S. I. (1998). Hometown Associations: Indigenous Knowledge and Development in Nigeria. Intermediate Technology Publications.
\bibitem{ifinedo2006} Ifinedo, P. (2006). Acceptance and continuance of e-government services in Nigeria.
\bibitem{iso2016} ISO 15489-1:2016. Information and documentation — Records management.
\bibitem{janowski2015} Janowski, T. (2015). Digital government evolution: From e-Government to T-government, Open Government and Smart Government. Government Information Quarterly.
\bibitem{langbein2010} Langbein, L., \& Knack, S. (2010). The World Bank’s Governance Indicators: What do they measure? Journal of Development Studies.
\bibitem{lemieux2016} Lemieux, V. L. (2016). Digital Records and the Evolution of the Digital Continuity Framework.
\bibitem{oecd2020} OECD. (2020). The OECD Digital Government Index (DGI): 2019 Results.
\bibitem{okafor1992} Okafor, F. U. (1992). Igbo Philosophy of Law. Fourth Dimension Publishers.
\bibitem{olson1965} Olson, M. (1965). The Logic of Collective Action. Harvard University Press.
\bibitem{oman2010} Oman, C. P., \& Arndt, C. (2010). Measuring governance. OECD Development Centre Policy Briefs.
\bibitem{ostrom1990} Ostrom, E. (1990). Governing the Commons. Cambridge University Press.
\bibitem{peffers2007} Peffers, K., et al. (2007). A design science research methodology for information systems research. Journal of Management Information Systems.
\bibitem{rowe1999} Rowe, G., \& Wright, G. (1999). The Delphi technique as a forecasting tool: Issues and analysis. International Journal of Forecasting.
\bibitem{schneider2019} Schneider, N. (2019). Decentralization: An exercise in discipline.
\bibitem{schouten2011} Schouten, M. A. (2011). The Social Audit Tool.
\bibitem{suchman1995} Suchman, M. C. (1995). Managing legitimacy: Strategic and institutional approaches. Academy of Management Review.
\bibitem{tyler2006} Tyler, T. R. (2006). Why people obey the law. Princeton University Press.
\bibitem{uduku2002} Uduku, O. (2002). The socio-economic impact of hometown associations. Africa.
\bibitem{un2022} United Nations. (2022). E-Government Survey 2022.
\bibitem{wb2021} World Bank. (2021). GovTech Maturity Index: The State of Public Sector Digital Transformation.

\end{thebibliography}
\end{document}
